% =============================
% Tractable Circuit Zoo - Query Support
% Auto-generated from database.json
% =============================
\documentclass[11pt]{article}
\usepackage[margin=1in]{geometry}
\usepackage{amsmath, amssymb, amsthm}
\usepackage{mathtools}
\usepackage{xparse}
\usepackage{hyperref}
\usepackage{natbib}
\hypersetup{colorlinks=true, linkcolor=blue, citecolor=blue, urlcolor=blue}



\NewDocumentCommand{\langref}{m g}{\textbf{#1\IfNoValueF{#2}{#2}}}
\NewDocumentCommand{\langfam}{m m g}{\textbf{#1$_{#2}$\IfNoValueF{#3}{#3}}}
\newcommand{\thislang}{\mathcal{L}}
\newcommand{\poly}{polynomial}
\newcommand{\nopolyunknownquasi}{not polynomial, quasipolynomial unknown}
\newcommand{\nopolyquasi}{not polynomial, quasipolynomial}
\newcommand{\unknownpolyquasi}{polynomial unknown, quasipolynomial}
\newcommand{\unknownboth}{unknown}
\newcommand{\noquasi}{not quasipolynomial}
\newcommand{\CO}{CO}
\newcommand{\VA}{VA}
\newcommand{\CE}{CE}
\newcommand{\IM}{IM}
\newcommand{\EQ}{EQ}
\newcommand{\SE}{SE}
\newcommand{\CT}{CT}
\newcommand{\ME}{ME}
\newcommand{\CD}{CD}
\newcommand{\FO}{FO}
\newcommand{\SFO}{SFO}
\newcommand{\NOTC}{$\neg$C}
\newcommand{\ANDC}{$\wedge$C}
\newcommand{\ANDBC}{$\wedge$BC}
\newcommand{\ORC}{$\vee$C}
\newcommand{\ORBC}{$\vee$BC}
\newcommand{\compilespoly}[2]{#1 compiles polynomially to #2}
\newcommand{\compilesquasi}[2]{#1 compiles quasipolynomially to #2}
\newcommand{\nocompilespoly}[2]{#1 does not compile polynomially to #2}
\newcommand{\nocompilesquasi}[2]{#1 does not compile quasipolynomially to #2}
\newcommand{\supportspoly}[2]{#1 supports #2 in polynomial time}
\newcommand{\supportsquasi}[2]{#1 supports #2 in quasipolynomial time}
\newcommand{\nosupportspoly}[2]{#2 is not supported by #1 in polynomial time}
\newcommand{\nosupportsquasi}[2]{#2 is not supported by #1 in quasipolynomial time}
\newcommand{\isin}[1]{#1}
\newcommand{\shortname}[1]{\textbf{Short name:} #1\par}
\newcommand{\fullname}[1]{\textbf{Full name:} #1\par}
\newcommand{\source}[1]{\textbf{Source:} #1\par}
\newcommand{\target}[1]{\textbf{Target:} #1\par}
\newcommand{\language}[1]{\textbf{Language:} #1\par}
\newcommand{\operation}[1]{\textbf{Operation:} #1\par}
\newcommand{\status}[1]{\textbf{Status:} #1\par}
\newcommand{\selector}[1]{\textbf{Selector:} #1\par}
\newcommand{\assuming}[1]{\textbf{Assuming:} #1\par}
\newcommand{\titlefield}[1]{\textbf{Title:} #1\par}
\let\titlemacro\title
\renewcommand{\title}[1]{\titlefield{#1}}
\newenvironment{description}{\par\noindent\ignorespaces}{\par}
\newenvironment{language}{}{}
\newenvironment{concept}{}{}
\newenvironment{succinctnessclaim}{}{}
\newenvironment{queryclaim}{}{}
\newenvironment{transformationclaim}{}{}
\newenvironment{batchclaim}{}{}
\titlemacro{Tractable Circuit Zoo: Query Support}
\date{\today}
\begin{document}
\maketitle

\begin{queryclaim}
\language{\langref{CNF}}
\operation{\CO}
\status{\nopolyunknownquasi}
\assuming{P \neq NP}
\begin{description}
Consistency checking for \langref{CNF} is NP-complete by the Cook-Levin theorem, the classic foundational result for SAT \citep{Cook_1971,Levin_1973}.
\end{description}
\end{queryclaim}

\begin{queryclaim}
\language{\langref{CNF}}
\operation{\IM}
\status{\poly}
\begin{description}
To check whether a consistent term $\alpha$ implies a \langref{CNF} representation $\Sigma$, it suffices to verify that each clause of $\Sigma$ shares a literal with $\alpha$. This can be done in polynomial time \citep{Darwiche_2002}.
\end{description}
\end{queryclaim}

\begin{queryclaim}
\language{\langref{CNF}}
\operation{\VA}
\status{\poly}
\begin{description}
A \langref{CNF} representation is valid iff every clause in it is valid (i.e., contains a complementary pair of literals). This can be checked in linear time \citep{Darwiche_2002}.
\end{description}
\end{queryclaim}

\begin{queryclaim}
\language{\langref{d-DNNF}}
\operation{\EQ}
\status{\poly}
\assuming{$P = \mathrm{BPP}$}
\begin{description}
\citet{Darwiche_Huang_2002} showed that equivalence of \langref{d-DNNF} representations can be tested in randomized polynomial time by reducing it to the problem of checking whether two \langref{d-DNNF} representations agree on a randomly chosen assignment. Under the derandomization hypothesis P = BPP, this yields a deterministic polynomial-time algorithm and hence \langref{d-DNNF} satisfies EQ.
\end{description}
\end{queryclaim}

\begin{queryclaim}
\language{\langref{d-DNNF}}
\operation{\CT}
\status{\poly}
\begin{description}
\langref{d-DNNF} satisfies CT: model counting can be performed in polynomial time by exploiting decomposability (multiplying counts at $\land$-nodes) and determinism (adding counts at $\lor$-nodes). This is the key tractability result motivating \langref{d-DNNF} as a compilation target \citep{Darwiche_2001b} \citep{Darwiche_2002}.
\end{description}
\end{queryclaim}

\begin{queryclaim}
\language{\langref{DNF}}
\operation{\VA}
\status{\nopolyunknownquasi}
\assuming{P \neq NP}
\begin{description}
Validity checking for \langref{DNF} is coNP-complete. A \langref{DNF} is valid iff every assignment satisfies at least one term, which requires checking that term coverage is complete \citep{Darwiche_2002}.
\end{description}
\end{queryclaim}

\begin{queryclaim}
\language{\langref{DNNF}}
\operation{\CO}
\status{\poly}
\begin{description}
Consistency of a \langref{DNNF} representation can be decided by a bottom-up traversal. A literal or constant node is handled directly; an $\lor$-node is consistent iff at least one child is consistent; and an $\land$-node is consistent iff all children are consistent. The last test is sound because decomposability ensures that conjuncts mention disjoint variable sets, so independently satisfiable children have compatible models. Hence consistency is decidable in linear time in the circuit size. \citep{Darwiche_2001a,Darwiche_2002}
\end{description}
\end{queryclaim}

\begin{queryclaim}
\language{\langref{FBDD}}
\operation{\CO}
\status{\poly}
\begin{description}
Consistency means existence of a satisfying root-to-1 path. For a general \langref{FBDD}, run graph reachability from the root to terminal 1 in linear time \citep{Gergov_1994,Bryant_1992}.
\end{description}
\end{queryclaim}

\begin{queryclaim}
\language{\langref{FBDD}}
\operation{\CT}
\status{\poly}
\begin{description}
Model counting is done by dynamic programming over the DAG representation of the free BDD, combining child counts at each decision node; total runtime is polynomial in graph size \citep{Gergov_1994,Bryant_1992}.
\end{description}
\end{queryclaim}

\begin{queryclaim}
\language{\langref{FBDD}}
\operation{\VA}
\status{\poly}
\begin{description}
Validity is equivalent to absence of any falsifying path. For a general \langref{FBDD}, check that terminal 0 is unreachable from the root (equivalently, all root-to-terminal paths evaluate to true); this is graph reachability in polynomial time \citep{Gergov_1994,Bryant_1992}.
\end{description}
\end{queryclaim}

\begin{queryclaim}
\language{\langref{FBDD}}
\operation{\EQ}
\status{\poly}
\assuming{$P = \mathrm{BPP}$}
\begin{description}
\citep{Blum_1980} evaluate both free Boolean graphs under random assignments over a finite field; if values differ, the graphs are definitely distinct, and if values match, they are probably equivalent. Repetition decreases error probability exponentially. Under derandomization assumptions such as P = BPP, this also yields deterministic polynomial-time equivalence testing for \langref{FBDD}.
\end{description}
\end{queryclaim}

\begin{queryclaim}
\language{\langref{IP}}
\operation{\CT}
\status{\nopolyunknownquasi}
\assuming{the polynomial hierarchy does not collapse}
\begin{description}
By duality with \langref{PI}: the negation of a formula in prime implicate form is in prime implicant (\langref{IP}) form. The number of models of $\neg\Sigma$ over $\text{Vars}(\Sigma)$ is $2^n$ minus the model count of $\Sigma$. Hence the \#P-hardness of counting for \langref{PI} transfers to \langref{IP} \citep{Darwiche_2002}.
\end{description}
\end{queryclaim}

\begin{queryclaim}
\language{\langref{IP}}
\operation{\IM}
\status{\poly}
\begin{description}
By definition, \langref{IP} is the set of prime implicants of a formula. A consistent term $\alpha$ implies $\Sigma$ iff $\alpha$ entails some prime implicant of $\Sigma$. This entailment check between terms runs in polynomial time \citep{Darwiche_2002}.
\end{description}
\end{queryclaim}

\begin{queryclaim}
\language{\langref{IP}}
\operation{\SE}
\status{\poly}
\begin{description}
\langref{IP} is a subset of \langref{DNF} and \langref{IP} satisfies IM, so checking sentential entailment reduces to checking implication between prime implicant forms. Two formulas are equivalent iff they share the same prime implicants (canonical form), giving SE \citep{Darwiche_2002}.
\end{description}
\end{queryclaim}

\begin{queryclaim}
\language{\langref{MODS}}
\operation{\CO}
\status{\poly}
\begin{description}
A representation $\Sigma$ is consistent iff it has at least one model. Given the set of models, we need only check nonemptiness of the set \citep{Darwiche_2002}.
\end{description}
\end{queryclaim}

\begin{queryclaim}
\language{\langref{MODS}}
\operation{\CT}
\status{\poly}
\begin{description}
We need only take the cardinality of the set of models (which can be done in linear time) \citep{Darwiche_2002}.
\end{description}
\end{queryclaim}

\begin{queryclaim}
\language{\langfam{OBDD}{<}}
\operation{\CO}
\status{\poly}
\begin{description}
Consistency can be checked in polynomial time by graph reachability: verify that there exists a path from the root to a terminal 1 node \citep{Bryant_1986}.
\end{description}
\end{queryclaim}

\begin{queryclaim}
\language{\langfam{OBDD}{<}}
\operation{\CT}
\status{\poly}
\begin{description}
Model counting on \langfam{OBDD}{<}{s} can be done in polynomial time by a single bottom-up pass \citep{Bryant_1986, Bryant_1992}.
\end{description}
\end{queryclaim}

\begin{queryclaim}
\language{\langfam{OBDD}{<}}
\operation{\EQ}
\status{\poly}
\begin{description}
After reducing both diagrams, \langfam{OBDD}{<}{s} are canonical: two functions are equivalent iff their reduced \langfam{OBDD}{<} diagrams are identical. This makes EQ a structural isomorphism check \citep{Bryant_1986, Bryant_1992, Meinel_Theobald_1998} (Theorem 8.11).
\end{description}
\end{queryclaim}

\begin{queryclaim}
\language{\langfam{OBDD}{<}}
\operation{\ME}
\status{\poly}
\begin{description}
Models of an \langfam{OBDD}{<} formula can be enumerated in time polynomial in the output size by traversing all root-to-1-sink paths. Each path corresponds to a (partial) assignment that can be expanded to full models. \citep{Bryant_1986}
\end{description}
\end{queryclaim}

\begin{queryclaim}
\language{\langfam{OBDD}{<}}
\operation{\VA}
\status{\poly}
\begin{description}
Validity can be checked in polynomial time by graph reachability: verify that no path from the root reaches the terminal 0 node, equivalently that every reachable terminal is 1 \citep{Bryant_1986}.
\end{description}
\end{queryclaim}

\begin{queryclaim}
\language{\langref{PI}}
\operation{\CE}
\status{\poly}
\begin{description}
By definition, \langref{PI} is the set of prime implicates of a formula. A clause $\gamma$ is entailed by $\Sigma$ iff $\gamma$ subsumes (is a weakening of) some prime implicate of $\Sigma$. This subsumption check runs in polynomial time \citep{Darwiche_2002}.
\end{description}
\end{queryclaim}

\begin{queryclaim}
\language{\langref{PI}}
\operation{\CT}
\status{\nopolyunknownquasi}
\assuming{the polynomial hierarchy does not collapse}
\begin{description}
The model counting problem for monotone Krom formulas (conjunctions of clauses with at most two positive literals) is \#P-complete \citep{Roth_1996}. Such formulas can be converted to prime implicates form in polynomial time, so \langref{PI} does not satisfy CT \citep{Darwiche_2002}.
\end{description}
\end{queryclaim}

\begin{queryclaim}
\language{\langref{PI}}
\operation{\SE}
\status{\poly}
\begin{description}
\langref{PI} is a subset of \langref{CNF} and \langref{PI} satisfies CE, so checking $\Sigma_1 \models \Sigma_2$ for \langref{CNF} $\Sigma_2$ reduces to checking each clause of $\Sigma_2$ against the prime implicates of $\Sigma_1$. Two formulas are equivalent iff they share the same prime implicates (canonical form), giving SE \citep{Darwiche_2002}.
\end{description}
\end{queryclaim}

\begin{queryclaim}
\language{\langref{OBDD}}
\operation{\CT}
\status{\poly}
\begin{description}
Any query concerning \langref{OBDD} is equivalent to the corresponding query concerning \langfam{OBDD}{<} when only one DAG is involved. Since \langfam{OBDD}{<} satisfies CO, VA and CT, so does \langref{OBDD} \citep{Darwiche_2002}. \citep{Bryant_1986}
\end{description}
\end{queryclaim}

\begin{queryclaim}
\language{\langref{OBDD}}
\operation{\EQ}
\status{\poly}
\begin{description}
\citet{Meinel_Theobald_1998} cite the polynomial-time equivalence algorithm of \citet{Fortune_1978} for free single-variable program schemes where one input tests predicates in one fixed order on all executions. An \langref{OBDD} is such a free read-once scheme: along every root-to-sink path, variables are tested according to its associated order. The algorithm compares the two diagrams by recursively comparing residual subprograms. When the ordered input next tests variable $x$, both $x=0$ and $x=1$ cofactors are compared against the corresponding restrictions of the other diagram; terminal disagreements reject, and repeated residual pairs are memoized. The number of residual pairs is polynomial in the input sizes, so equivalence of two \langref{OBDD} representations is decidable in polynomial time \citep{Darwiche_2002}.
\end{description}
\end{queryclaim}

\begin{queryclaim}
\language{\langref{OBDD}}
\operation{\SE}
\status{\nopolyunknownquasi}
\assuming{P \neq NP}
\begin{description}
Sentential entailment for \langref{OBDD} is coNP-complete. \citet[Lemma 8.14]{Meinel_Theobald_1998} prove that checking consistency of two \langfam{OBDD}{<}{s} using different variable orders is NP-complete. Since \langref{OBDD} supports negation, consistency of $\alpha \land \beta$ is equivalent to $\alpha \not\models \neg\beta$, so polynomial-time SE for arbitrary \langref{OBDD}{s} would imply polynomial-time consistency for such pairs \citep{Darwiche_2002}.
\end{description}
\end{queryclaim}
\bibliographystyle{plainnat}
\bibliography{refs}
\end{document}
